\chapter{Lógica Secuencial}

\section{Lógica combinacional vs. lógica secuencial}

En lógica combinacional, las salidas dependen únicamente del valor actual de las entradas. En lógica \textbf{secuencial}, las salidas dependen tanto de las entradas actuales como del \textbf{estado anterior} del circuito. Esto significa que el circuito tiene memoria: recuerda lo que ocurrió antes.

El bloque fundamental de la lógica secuencial es el \textbf{dispositivo biestable}: un circuito que puede permanecer indefinidamente en uno de dos estados estables (SET o RESET), y que solo cambia de estado cuando se lo indica una señal de entrada.

Los dispositivos biestables se dividen en dos categorías:

\begin{itemize}
    \item \textbf{Latches:} cambian de estado de forma asíncrona, es decir, responden directamente al nivel de sus entradas en cualquier momento.
    \item \textbf{Flip-flops:} cambian de estado de forma síncrona, es decir, solo en el momento en que ocurre un flanco de la señal de reloj.
\end{itemize}

\section{Latches}

\subsection{Latch SR (Set-Reset)}

Es el tipo más simple de dispositivo biestable. Tiene dos entradas (S y R) y dos salidas complementarias (Q y $\overline{Q}$).

Puede implementarse con dos puertas NOR (entradas activas a nivel alto) o con dos puertas NAND (entradas activas a nivel bajo).

\begin{figure}[H]
\centering
% ============================================================
%  Latch SR — implementación con compuertas (figura 7.1 Floyd)
%   Izquierda:  con NOR  -> entradas S, R activas en ALTO
%   Derecha:    con NAND -> entradas Sbar, Rbar activas en BAJO
%  Las compuertas están acopladas en CRUZ: la salida de cada una
%  realimenta una entrada de la otra (el lazo de realimentación se
%  dibuja rodeando el par de compuertas).
%
%  LÓGICA (importante): en el latch NOR, S debe poner Q=1. Como el NOR
%  invierte, S entra en la compuerta cuya salida es Qbar (abajo) y R en
%  la de salida Q (arriba). Por eso R queda arriba y S abajo.
%  (En el NAND activo-bajo, Sbar arriba sí produce Q arriba.)
% ============================================================
%  CÓMO AJUSTAR:
%   - Separación vertical de compuertas: el y=±1.2 de cada \node
%   - Separación NOR/NAND:               el shift={(8,0)} del 2º scope
%   - Tamaño del lazo de realimentación: los ++(0,±4) y ++(-5.x,0)
%   NOTA: el único cruce de cable SIN punto (realim. Qbar sobre la
%         línea de Q) NO conecta; el punto (\circ) marca la derivación.
% ============================================================
\begin{circuitikz}[transform shape, scale=0.8]

  % ===================== Latch con NOR (activo en alto) =====================
  \node[ieeestd nor port] (a) at (0, 1.2) {};    % salida Q
  \node[ieeestd nor port] (b) at (0,-1.2) {};    % salida Qbar
  \draw (a.in 1) -- ++(-1.2,0)node[left]{\small $R$};
  \draw (b.in 2) -- ++(-1.2,0)node[left]{\small $S$};
  \draw (a.out) -- ++(0.2,0) coordinate(qa) -- ++(0.6,0) node[right]{\small $Q$}            ;
  \draw (b.out) -- ++(0.2,0) coordinate(qb) -- ++(0.6,0) node[right]{\small $\overline{Q}$} ;
  \node[circ] at (qa){};   
  \node[circ] at (qb){};
  \draw (qa) -- ++(0,-0.5) --++(-2.5,-1) |- (b.in 1);
  \draw (qb) -- ++(0, 0.5) --++(-2.5, 1) |- (a.in 2);


  % ===================== Latch con NAND (activo en bajo) =====================
  \begin{scope}[shift={(8,0)}]
    \node[ieeestd nand port] (c) at (0, 1.2) {};   % salida Q
    \node[ieeestd nand port] (d) at (0,-1.2) {};   % salida Qbar
    \draw (c.in 1) -- ++(-1.2,0)node[left]{\small $\overline{S}$};
    \draw (d.in 2) -- ++(-1.2,0)node[left]{\small $\overline{R}$};
    \draw (c.out) -- ++(0.2,0) coordinate(qc) -- ++(0.6,0) node[right]{\small $Q$}            ;
    \draw (d.out) -- ++(0.2,0) coordinate(qd) -- ++(0.6,0) node[right]{\small $\overline{Q}$} ;
    \node[circ] at (qc){};   
    \node[circ] at (qd){};
    \draw (qc) -- ++(0,-0.5) --++(-2.5,-1) |- (d.in 1);
    \draw (qd) -- ++(0, 0.5) --++(-2.5, 1) |- (c.in 2);
  \end{scope}

\end{circuitikz}

\caption{Latch SR: implementación con puertas NOR (izquierda, entradas activas alto)
  y puertas NAND (derecha, entradas activas bajo).
  Los cruces de cables sin punto indican que los conductores no se conectan.}
\end{figure}

\tikzCircuito{cap1/latch-SR-simbolo}{Símbolo lógico del latch SR: representación en bloque con entradas S, R y salidas Q, $\overline{Q}$ (figura 7.4 Floyd)}

\textbf{Tabla de verdad --- Latch SR (puertas NOR, entradas activas alto):}

\begin{table}[H]
\centering
\begin{tabular}{ccccp{6cm}}
\toprule
$S$ & $R$ & $Q$ & $\overline{Q}$ & Condición \\
\midrule
0 & 0 & $Q_0$ & $\overline{Q_0}$ & No cambia (retiene estado) \\
1 & 0 & 1 & 0 & SET (Q pasa a 1) \\
0 & 1 & 0 & 1 & RESET (Q pasa a 0) \\
1 & 1 & ? & ? & \textbf{No válido} \\
\bottomrule
\end{tabular}
\end{table}

La condición no válida ocurre cuando ambas entradas se activan simultáneamente: fuerza a ambas salidas al mismo nivel, violando la condición de complementariedad. Si luego se desactivan al mismo tiempo, el estado final es impredecible.

\bannerRevision
\begin{figure}[H]
\centering
\begin{tikztimingtable}[
  timing/slope=0,
  timing/coldist=2pt,
  timing/rowdist=3.8ex,
  timing/name/.style={font=\small\sffamily, right},
  xscale=0.9
]
  $S$            & 2L 2H 8L \\
  $R$            & 8L 2H 2L \\
  $Q$            & 2X 6H 4L \\
  $\overline{Q}$ & 2X 6L 4H \\
  \extracode
    % líneas de referencia en las transiciones
    \begin{pgfonlayer}{background}
      \draw[lightgray, dashed] (2,0.4) -- (2,-5.2);
      \draw[lightgray, dashed] (4,0.4) -- (4,-5.2);
      \draw[lightgray, dashed] (8,0.4) -- (8,-5.2);
      \draw[lightgray, dashed] (10,0.4) -- (10,-5.2);
    \end{pgfonlayer}
    % etiquetas de fase (encima de la primera señal)
    \node[font=\footnotesize\sffamily,above] at (1, 0.4)  {Hold};
    \node[font=\footnotesize\sffamily,above] at (3, 0.4)  {SET};
    \node[font=\footnotesize\sffamily,above] at (6, 0.4)  {Hold};
    \node[font=\footnotesize\sffamily,above] at (9, 0.4)  {RESET};
    \node[font=\footnotesize\sffamily,above] at (11,0.4)  {Hold};
\end{tikztimingtable}

\caption{Diagrama de tiempos del latch SR (puertas NOR, entradas activas alto).
  La franja \textit{Hold} indica que el estado se retiene sin cambios.}
\end{figure}

\textbf{Aplicación --- Eliminación de rebote de pulsadores:}

\subsubsection*{¿Qué es el rebote de los contactos?}

Cuando se presiona un pulsador o se acciona un interruptor mecánico, uno esperaría que la señal eléctrica pasara limpiamente de 0 a 1 en un único instante. Sin embargo, los contactos metálicos son físicamente elásticos: al cerrarse, rebotan varias veces antes de quedar en contacto estable. Lo mismo ocurre al abrir el circuito.

Este fenómeno se llama \textbf{rebote de contactos} (\textit{contact bounce} o \textit{switch bounce}). La duración del rebote depende del tipo de pulsador y de su estado mecánico, pero típicamente está entre \textbf{1\,ms y 20\,ms}.

\tikzCircuito{cap1/bounce-signal}{Señal ideal (flanco limpio) vs.\ señal real con rebote (múltiples pulsos breves en la zona de transición)}

\subsubsection*{¿Qué ve el circuito digital durante el rebote?}

Para la electrónica digital no existe el concepto de ``está rebotando'': cada transición de tensión es leída como un cambio de estado válido. Así, un único pulso de dedo puede generar \textbf{5, 10 o más pulsos eléctricos} en los primeros milisegundos. La señal real se ve como una ráfaga de pulsos muy cortos, seguida por un nivel estable. Desde el punto de vista del circuito, esos pulsos son completamente válidos e indistinguibles de pulsaciones reales del usuario.

\subsubsection*{Problemas que genera el rebote}

\begin{table}[H]
\centering
\begin{tabular}{p{5.5cm} p{7cm}}
\toprule
\textbf{Situación} & \textbf{Consecuencia del rebote} \\
\midrule
Contador incrementado por botón & Un solo toque incrementa el contador en 3, 7 u otro número imprevisible \\
\addlinespace
Interrupción por flanco en microcontrolador & La ISR se ejecuta múltiples veces por una única pulsación \\
\addlinespace
Latch o flip-flop activado por botón & El estado puede cambiar varias veces antes de estabilizarse \\
\addlinespace
Sistema de navegación de menú & Un toque avanza varios ítems en lugar de uno \\
\bottomrule
\end{tabular}
\caption{Problemas típicos causados por el rebote de contactos en sistemas digitales.}
\end{table}

El rebote es especialmente problemático con señales de reloj o con pines de interrupción, donde el sistema es sensible a cada flanco individualmente.

\subsubsection*{Solución con latch SR}

Un latch SR puede eliminar el rebote mecánico usando un interruptor SPDT (\textit{Single Pole Double Throw}) conectado entre masa y las entradas $\overline{S}$ y $\overline{R}$ del latch, con resistencias pull-up a Vcc.

Al cambiar la posición del interruptor, el primer contacto activa el latch de forma definitiva. Cualquier rebote posterior aplica la misma entrada (S o R) que ya fue procesada, por lo que \textbf{el estado del latch no vuelve a cambiar}. La salida Q presenta una transición limpia y única.

\tikzCircuito{cap1/debounce-latch}{Circuito eliminador de rebote de pulsadores con latch SR (figura 7.6 Floyd)}

\tikzCircuito{cap1/debounce-timing}{Diagrama de tiempos del circuito antirrebote: señal con rebote en la entrada vs.\ salida Q limpia del latch (figura 7.7 Floyd)}

\subsection{Latch SR con habilitación}

Agrega una entrada de habilitación (EN). El latch solo puede cambiar de estado cuando EN está activo. Permite controlar exactamente en qué momento el latch responde a sus entradas.

\tikzCircuito{cap1/latch-SR-EN}{Diagrama y símbolo del latch SR con habilitación (EN) (figura 7.8 Floyd)}

\textbf{Tabla de verdad --- Latch SR con habilitación (puertas NOR, entradas activas alto):}

\begin{table}[H]
\centering
\begin{tabular}{ccccp{6.5cm}}
\toprule
$S$ & $R$ & EN & $Q$ & Condición \\
\midrule
0 & 0 & 1 & $Q_0$ & No cambia \\
1 & 0 & 1 & 1     & SET (Q pasa a 1) \\
0 & 1 & 1 & 0     & RESET (Q pasa a 0) \\
1 & 1 & 1 & ?     & \textbf{No válido} \\
X & X & 0 & $Q_0$ & No cambia (EN inactivo, S y R ignoradas) \\
\bottomrule
\end{tabular}
\end{table}

\bannerRevision
\begin{figure}[H]
\centering
\begin{tikztimingtable}[
  timing/slope=0,
  timing/coldist=2pt,
  timing/rowdist=3.8ex,
  timing/name/.style={font=\small\sffamily, right},
  xscale=0.9
]
  EN  & 4H 4L 4H \\
  $S$ & 2L 2H 8L \\
  $R$ & 5L 5H 2L \\
  $Q$ & 2L 6H 4L \\
  \extracode
    \begin{pgfonlayer}{background}
      \draw[lightgray, dashed] (2, 0.4) -- (2,-5.2);
      \draw[lightgray, dashed] (4, 0.4) -- (4,-5.2);
      \draw[lightgray, dashed] (5, 0.4) -- (5,-5.2);
      \draw[lightgray, dashed] (8, 0.4) -- (8,-5.2);
      \draw[lightgray, dashed] (10,0.4) -- (10,-5.2);
    \end{pgfonlayer}
    \node[font=\footnotesize\sffamily,above] at (1,  0.4) {Hold};
    \node[font=\footnotesize\sffamily,above] at (3,  0.4) {SET};
    \node[font=\footnotesize\sffamily,above] at (6.5,0.4) {EN=0 (bloqueado)};
    \node[font=\footnotesize\sffamily,above] at (9,  0.4) {RESET};
    \node[font=\footnotesize\sffamily,above] at (11, 0.4) {Hold};
    \node[font=\footnotesize\itshape, anchor=west] at (0,-6.3)
      {R se activa en $t{=}5$ pero es ignorado hasta que EN vuelve a 1 (en $t{=}8$).};
\end{tikztimingtable}

\caption{Diagrama de tiempos del latch SR con habilitación. EN=0 bloquea las entradas:
  aunque R se active mientras EN=0, Q no cambia. El RESET recién ocurre
  cuando EN regresa a 1 (y R sigue activo).}
\end{figure}

\subsection{Latch D con habilitación}

Tiene una única entrada de datos (D) además de la habilitación (EN). Elimina la condición no válida del latch SR: cuando EN está activo, la salida Q sigue al valor de D. Cuando EN se desactiva, el latch retiene el último valor de D.

\tikzCircuito{cap1/latch-D}{Diagrama y símbolo del latch D con habilitación (figura 7.10 Floyd)}

\begin{table}[H]
\centering
\begin{tabular}{ccp{2cm}l}
\toprule
D & EN & Q & Condición \\
\midrule
0 & 1 & 0 & RESET \\
1 & 1 & 1 & SET \\
X & 0 & $Q_0$ & No cambia \\
\bottomrule
\end{tabular}
\end{table}

\bannerRevision
\begin{figure}[H]
\centering
\begin{tikztimingtable}[
  timing/slope=0,
  timing/coldist=2pt,
  timing/rowdist=3.8ex,
  timing/name/.style={font=\small\sffamily, right},
  xscale=0.9
]
  EN  & 4H 4L 4H \\
  $D$ & 2L 2H 2H 2L 2L 2H \\
  $Q$ & 2L 6H 2L 2H \\
  \extracode
    \begin{pgfonlayer}{background}
      % Zona EN=0 sombreada
      \fill[yellow!12] (4,-0.5) rectangle (8,-4.2);
      \draw[lightgray, dashed] (2, 0.4) -- (2,-4.2);
      \draw[lightgray, dashed] (4, 0.4) -- (4,-4.2);
      \draw[lightgray, dashed] (6, 0.4) -- (6,-4.2);
      \draw[lightgray, dashed] (8, 0.4) -- (8,-4.2);
      \draw[lightgray, dashed] (10,0.4) -- (10,-4.2);
    \end{pgfonlayer}
    \node[font=\footnotesize\sffamily,above] at (1, 0.4) {Q=D};
    \node[font=\footnotesize\sffamily,above] at (3, 0.4) {Q=D};
    \node[font=\footnotesize\sffamily,above] at (6, 0.4) {EN=0: Q retiene};
    \node[font=\footnotesize\sffamily,above] at (9, 0.4) {Q=D};
    \node[font=\footnotesize\sffamily,above] at (11,0.4) {Q=D};
    % Destacar cambio de D ignorado
    \node[font=\tiny\itshape, teal!70!black] at (7,-4.8)
      {D cambia a 0 pero Q no cambia (EN=0)};
\end{tikztimingtable}

\caption{Diagrama de tiempos del latch D. Con EN=1 la salida es \textit{transparente}:
  Q sigue inmediatamente a D. Con EN=0 el latch queda \textit{opaco}: retiene el
  último valor de D aunque la entrada cambie (zona sombreada; ver $t{=}6$).}
\end{figure}

\section{Flip-flops}

La diferencia fundamental entre un latch y un flip-flop es el mecanismo de disparo. Un flip-flop \textbf{disparado por flanco} solo cambia de estado en el instante exacto del flanco (de subida o de bajada) de la señal de reloj. Fuera de ese instante, es insensible a los cambios en sus entradas.

En los símbolos lógicos, el flanco activo se indica con un \textbf{triángulo} en la entrada de reloj (indicador de entrada dinámica). Si hay un círculo adicional, el disparo es por flanco negativo (bajada); sin círculo, por flanco positivo (subida).

\bannerRevision
\begin{figure}[H]
\centering
% ============================================================
%  Símbolos de flip-flops disparados por flanco
%  - FF-D flanco ↑ : cuña (triángulo) en la entrada de reloj.
%  - FF-D flanco ↓ : cuña + BURBUJA (círculo de negación) en el reloj.
%  - FF-JK flanco ↑: J, CLK, K.
%  Los FF de circuitikz se definen con "flipflop def" (manual §4.23):
%    tN = texto del pin N, cN=1 dibuja la cuña de reloj,
%    nN=1 dibuja la burbuja de negación en el pin N.
%  Pines: 1,2,3 a la izquierda (arriba→abajo); 6,5,4 a la derecha.
% ============================================================
%  CÓMO AJUSTAR:
%   - Separación entre símbolos: la x de cada \node (0, 5.4, 10.8)
%   - Largo de los cables de entrada/salida: el 0.6 de los ++(...)
%   - Tamaño general: scale= en \begin{circuitikz}
% ============================================================
\begin{circuitikz}[scale=0.85, transform shape,
  flipflop Dn/.style={flipflop,
    flipflop def={t1=D, t6=Q, c3=1, n3=1, t4={\ctikztextnot{Q}}}},
]

  % --- FF-D flanco positivo ---
  \node[flipflop D] (ffd) at (0,0) {};
  \draw (ffd.pin 1) -- ++(-0.6,0) node[left]{\small D};
  \draw (ffd.pin 3) -- ++(-0.6,0) node[left]{\small CLK};
  \draw (ffd.pin 6) -- ++(0.6,0) node[right]{\small $Q$};
  \draw (ffd.pin 4) -- ++(0.6,0) node[right]{\small $\overline{Q}$};
  \node[below=0.35 of ffd.south, font=\small] {Flanco $\uparrow$};

  % --- FF-D flanco negativo (burbuja en CLK) ---
  \node[flipflop Dn] (ffdn) at (5.4,0) {};
  \draw (ffdn.pin 1) -- ++(-0.6,0) node[left]{\small D};
  \draw (ffdn.pin 3) -- ++(-0.6,0) node[left]{\small CLK};
  \draw (ffdn.pin 6) -- ++(0.6,0) node[right]{\small $Q$};
  \draw (ffdn.pin 4) -- ++(0.6,0) node[right]{\small $\overline{Q}$};
  \node[below=0.35 of ffdn.south, font=\small] {Flanco $\downarrow$};

  % --- FF-JK flanco positivo ---
  \node[flipflop JK] (ffjk) at (10.8,0) {};
  \draw (ffjk.pin 1) -- ++(-0.6,0) node[left]{\small J};
  \draw (ffjk.pin 2) -- ++(-0.6,0) node[left]{\small CLK};
  \draw (ffjk.pin 3) -- ++(-0.6,0) node[left]{\small K};
  \draw (ffjk.pin 6) -- ++(0.6,0) node[right]{\small $Q$};
  \draw (ffjk.pin 4) -- ++(0.6,0) node[right]{\small $\overline{Q}$};
  \node[below=0.35 of ffjk.south, font=\small] {Flanco $\uparrow$};

\end{circuitikz}

\caption{Símbolos de flip-flops disparados por flanco. El triángulo en CLK indica
  disparo por flanco; el círculo adicional (centro) indica flanco negativo.}
\end{figure}

\subsection{Flip-flop SR disparado por flanco}

Similar al latch SR pero sincronizado con el reloj. Las entradas S y R solo son evaluadas en el flanco activo del reloj.

\begin{table}[H]
\centering
\begin{tabular}{cccc}
\toprule
S & R & Q siguiente & Condición \\
\midrule
0 & 0 & $Q_0$ & No cambia \\
1 & 0 & 1 & SET \\
0 & 1 & 0 & RESET \\
1 & 1 & ? & No válido \\
\bottomrule
\end{tabular}
\end{table}

\bannerRevision
\begin{figure}[H]
\centering
% ============================================================
%  Flip-flop SR disparado por flanco ascendente: diagrama de tiempos
%  - Las entradas S y R sólo se miran en los flancos ↑ de CLK
%    (t = 1, 3, 5, 7, 9). Entre flancos, Q no cambia.
%  - S y R cambian entre flancos (t par) para no violar setup/hold.
%  Coordenadas: ver tikz/estilos.tex (1 carácter = 1 unidad en x).
% ============================================================
\begin{tikztimingtable}[timing/aed]
  CLK & LHLHLHLHLH \\
  $S$ & 2L 2H 6L \\
  $R$ & 6L 2H 2L \\
  $Q$ & 3L 4H 3L \\
  \extracode
    \begin{pgfonlayer}{background}
      \vertlines{1,3,...,9}
    \end{pgfonlayer}
    \foreach \x/\op in {1/Hold, 3/SET, 5/Hold, 7/RESET, 9/Hold}
      \node[fase] at (\x,1.5) {\op};
\end{tikztimingtable}

\caption{Diagrama de tiempos del flip-flop SR disparado por flanco positivo.
  S y R solo son evaluadas en el flanco de subida del reloj; fuera de ese instante
  los cambios en las entradas no tienen efecto.}
\end{figure}

\subsection{Flip-flop D disparado por flanco}

La salida Q toma el valor de la entrada D en cada flanco activo del reloj. Es el más simple y el más utilizado en la práctica.

\begin{table}[H]
\centering
\begin{tabular}{cc}
\toprule
D & Q siguiente \\
\midrule
0 & 0 \\
1 & 1 \\
\bottomrule
\end{tabular}
\end{table}

\bannerRevision
\begin{figure}[H]
\centering
% ============================================================
%  Flip-flop D disparado por flanco ascendente: diagrama de tiempos
%  - Q copia el valor de D SÓLO en los flancos ↑ de CLK (t = 1,3,5,7,9).
%  - D cambia siempre entre flancos (t par) para no violar setup/hold.
%  Coordenadas: ver tikz/estilos.tex (1 carácter = 1 unidad en x).
% ============================================================
\begin{tikztimingtable}[timing/aed]
  CLK & LHLHLHLHLH \\
  $D$ & 4H 4L 2H   \\
  $Q$ & X 4H 4L H  \\
  \extracode
    \begin{pgfonlayer}{background}
      \vertlines{1,3,...,9}
    \end{pgfonlayer}
    % marca de captura en cada flanco activo
    \foreach \x in {1,3,...,9}
      \node[fase] at (\x,1.5) {$\uparrow$};
    \node[nota] at (0,{-2*\rowdist-1.1})
      {Q copia el valor de D sólo en cada flanco $\uparrow$ de CLK.};
\end{tikztimingtable}

\caption{Diagrama de tiempos del flip-flop D disparado por flanco positivo.
  Cada flanco de subida del reloj transfiere el valor de D a la salida Q.}
\end{figure}

\subsection{Flip-flop JK disparado por flanco}

El flip-flop JK resuelve la condición no válida del SR reemplazando el estado indeterminado por una operación de \textbf{basculación (toggle)}: cuando J = K = 1, la salida cambia de estado en cada flanco de reloj.

\begin{table}[H]
\centering
\begin{tabular}{cccc}
\toprule
J & K & Q siguiente & Condición \\
\midrule
0 & 0 & $Q_0$ & No cambia \\
1 & 0 & 1 & SET \\
0 & 1 & 0 & RESET \\
1 & 1 & $\overline{Q_0}$ & Bascula (Toggle) \\
\bottomrule
\end{tabular}
\end{table}

\tikzCircuito{cap1/flipflop-JK}{Símbolo del flip-flop JK disparado por flanco (figura 7.22 Floyd)}

\subsection{Flip-flop T (Toggle)}

Es un caso especial del JK con ambas entradas unidas (J = K = T). Cuando T = 1, la salida bascula en cada flanco de reloj. Cuando T = 0, la salida no cambia.

\begin{table}[H]
\centering
\begin{tabular}{cc}
\toprule
T & Q siguiente \\
\midrule
0 & $Q_0$ \\
1 & $\overline{Q_0}$ \\
\bottomrule
\end{tabular}
\end{table}

\bannerRevision
\begin{figure}[H]
\centering
% ============================================================
%  Flip-flop T disparado por flanco ascendente: diagrama de tiempos
%  - T=1 en el flanco ↑: Q se invierte (Toggle).
%  - T=0 en el flanco ↑: Q se mantiene (Hold).
%  - Flancos activos en t = 1, 3, 5, 7, 9.
%  Coordenadas: ver tikz/estilos.tex (1 carácter = 1 unidad en x).
% ============================================================
\begin{tikztimingtable}[timing/aed]
  CLK & LHLHLHLHLH \\
  $T$ & 2L 6H 2L \\
  $Q$ & 3L 2H 2L 3H \\
  \extracode
    \begin{pgfonlayer}{background}
      \vertlines{1,3,...,9}
    \end{pgfonlayer}
    \foreach \x/\op in {1/Hold, 3/Toggle, 5/Toggle, 7/Toggle, 9/Hold}
      \node[fase] at (\x,1.5) {\op};
\end{tikztimingtable}

\caption{Diagrama de tiempos del flip-flop T. Con T=0 la salida no cambia
  en el flanco activo; con T=1 la salida bascula en cada flanco de subida,
  dividiendo la frecuencia del reloj por 2.}
\end{figure}

El flip-flop T es la base de los contadores binarios.

\subsection{Entradas asíncronas: Preset y Clear}

La mayoría de los flip-flops tienen entradas asíncronas que actúan \textbf{independientemente del reloj}:

\begin{itemize}
    \item \textbf{Preset (PRE):} fuerza la salida Q a 1.
    \item \textbf{Clear (CLR):} fuerza la salida Q a 0.
\end{itemize}

Son útiles para inicializar el sistema a un estado conocido al encender, o para hacer resets de emergencia.

\subsection{Características de funcionamiento}

Algunas características importantes de los flip-flops al trabajar con señales reales:

\textbf{Retardo de propagación ($t_{pd}$):} tiempo que transcurre desde el flanco activo del reloj hasta que la salida estabiliza su nuevo valor. Limita la frecuencia máxima de operación.

\textbf{Tiempo de establecimiento ($t_{setup}$):} tiempo mínimo que la entrada debe estar estable \textbf{antes} del flanco activo del reloj para garantizar una lectura correcta.

\textbf{Tiempo de mantenimiento ($t_{hold}$):} tiempo mínimo que la entrada debe permanecer estable \textbf{después} del flanco activo del reloj.

\textbf{Frecuencia máxima ($f_{max}$):} frecuencia de reloj máxima a la que el flip-flop puede operar correctamente.

\bannerRevision
\begin{figure}[H]
\centering
\begin{tikzpicture}[
  >=Stealth,
  sig/.style={thick},
  dim/.style={<->,>=Stealth,teal!70!black,thick},
  anno/.style={font=\small\sffamily,teal!80!black},
  lbl/.style={font=\footnotesize\sffamily},
]
  % --- Parámetros de tiempo (unidades arbitrarias) ---
  \def\Tsetup{1.5}   % duración t_setup
  \def\Tclkedge{4}   % posición del flanco de CLK
  \def\Thold{1.5}    % duración t_hold
  \def\Tpd{1.8}      % retardo de propagación (CLK→Q válido)
  \def\Tend{8}

  % --- CLK ---
  \draw[sig] (0,2) -- (\Tclkedge-0.1,2)
             -- (\Tclkedge,3) -- (\Tend,3);   % flanco de subida en \Tclkedge
  \node[lbl,left] at (0,2.5) {CLK};

  % --- D (estable alrededor del flanco) ---
  % D=H antes y después del flanco
  \draw[sig] (0,0.8) -- (1.0,0.8)                  % inicio con dato previo (L)
             -- (1.0,0) -- (\Tclkedge-\Tsetup,0)   % transición de D
             -- (\Tclkedge-\Tsetup,0.8)             % D estable = H desde t_setup antes
             -- (\Tclkedge+\Thold,0.8)          % D estable hasta después de t_hold
             -- (\Tclkedge+\Thold,0)            % D puede cambiar
             -- (\Tend,0);
  \node[lbl,left] at (0,0.4) {D};

  % zona estable de D (relleno, detrás de las señales)
  \begin{pgfonlayer}{background}
    \fill[teal!12] (\Tclkedge-\Tsetup,0) rectangle (\Tclkedge+\Thold,0.8);
  \end{pgfonlayer}

  % --- Q ---
  \draw[sig] (0,-1.2) -- (\Tclkedge+\Tpd,-1.2)  % Q mantiene valor anterior
             -- (\Tclkedge+\Tpd,-0.4)             % Q cambia al nuevo valor
             -- (\Tend,-0.4);
  \node[lbl,left] at (0,-0.8) {Q};
  % zona de propagación (Q no válido)
  \begin{pgfonlayer}{background}
    \fill[orange!15] (\Tclkedge,-1.2) rectangle (\Tclkedge+\Tpd,-0.4);
  \end{pgfonlayer}
  \draw[orange!60,dashed] (\Tclkedge,-1.3) -- (\Tclkedge+\Tpd,-1.3);

  % --- Línea del flanco de CLK ---
  \draw[gray,dashed,thin] (\Tclkedge,-1.5) -- (\Tclkedge,3.5);

  % --- Cotas ---
  % t_setup
  \draw[dim] (\Tclkedge-\Tsetup,1.5) -- (\Tclkedge,1.5)
    node[midway,above,anno]{$t_{setup}$};

  % t_hold
  \draw[dim] (\Tclkedge,1.5) -- (\Tclkedge+\Thold,1.5)
    node[midway,above,anno]{$t_{hold}$};

  % t_pd
  \draw[dim] (\Tclkedge,-1.8) -- (\Tclkedge+\Tpd,-1.8)
    node[midway,below,anno]{$t_{pd}$};

  % flanco de subida
  \draw[->,gray] (\Tclkedge+0.3,2.4) node[right,lbl]{flanco activo} -- (\Tclkedge+0.05,2.8);

\end{tikzpicture}

\caption{Parámetros temporales de un flip-flop: $t_{setup}$ (D debe estar estable antes del flanco),
  $t_{hold}$ (D debe mantenerse estable después del flanco),
  $t_{pd}$ (tiempo desde el flanco hasta que Q es válido).
  La zona sombreada en D indica la ventana de estabilidad requerida.}
\end{figure}

\section{Registros de desplazamiento}

Un registro de desplazamiento es un circuito formado por flip-flops en cascada que cumple dos funciones: \textbf{almacenamiento} y \textbf{desplazamiento} de datos. Cada flip-flop representa un bit; por lo tanto, un registro de N etapas almacena N bits.

Los datos pueden entrar y salir del registro de distintas formas:

\imagenPendiente{los 7 modos básicos de movimiento de datos en un registro (figura 9.2 Floyd)}

\subsection{Entrada serie / Salida serie (SISO)}

Los bits entran uno a uno por un extremo y salen uno a uno por el otro. En cada pulso de reloj, todos los bits se desplazan una posición. Se necesitan N pulsos para cargar N bits y otros N para extraerlos.

\tikzCircuito{cap1/registro-SISO}{Registro de desplazamiento SISO de 4 bits: flip-flops D en cascada (figura 9.3 Floyd)}

\begin{notabox}
\textbf{CI típico:} \textbf{74HC164} --- registro SIPO de 8 bits (serie in / paralelo out). También puede usarse como SISO tomando solo la salida del último flip-flop.
\end{notabox}

\subsection{Entrada serie / Salida paralelo (SIPO)}

Los bits entran en serie pero se leen simultáneamente de todas las salidas Q de cada etapa. Una vez cargados los N bits, están disponibles todos al mismo tiempo en las salidas paralelo. Se usa para convertir datos serie a paralelo, por ejemplo para controlar múltiples LEDs o segmentos usando solo unos pocos pines del microcontrolador.

\begin{notabox}
\textbf{CI típicos:}
\begin{itemize}
    \item \textbf{74HC595} --- registro SIPO de 8 bits con latch de salida separado y salida en colector abierto para cascada. Es el más utilizado en proyectos con microcontroladores para expandir salidas digitales.
    \item \textbf{74HC164} --- registro SIPO de 8 bits sin latch de salida (las salidas cambian mientras se desplazan los datos).
\end{itemize}
\end{notabox}

\imagenPendiente{diagrama de pines y conexión del 74HC595 en cascada con microcontrolador}

\subsection{Entrada paralelo / Salida serie (PISO)}

Los N bits se cargan simultáneamente en las etapas (carga paralelo) y luego se extraen uno a uno en serie. Se usa para convertir datos paralelo a serie, por ejemplo para leer el estado de múltiples entradas digitales usando pocos pines del microcontrolador.

\begin{notabox}
\textbf{CI típicos:}
\begin{itemize}
    \item \textbf{74HC165} --- registro PISO de 8 bits con carga paralela asíncrona. Muy usado para leer matrices de pulsadores o DIP switches desde un microcontrolador por SPI.
    \item \textbf{74LS165} --- versión TTL del mismo.
\end{itemize}
\end{notabox}

\subsection{Entrada paralelo / Salida paralelo (PIPO)}

Los datos se cargan y se leen en paralelo. Funciona esencialmente como un banco de flip-flops D con habilitación. Se usa principalmente como registro de almacenamiento temporal o buffer de datos.

\begin{notabox}
\textbf{CI típico:} \textbf{74HC195} --- registro PIPO de 4 bits con entrada serie y paralelo y salida paralelo.
\end{notabox}

\subsection{Registros bidireccionales}

Pueden desplazar datos tanto hacia la izquierda como hacia la derecha, según el valor de una entrada de control de dirección. Son los más versátiles ya que combinan múltiples modos de operación.

\begin{notabox}
\textbf{CI típicos:}
\begin{itemize}
    \item \textbf{74HC194} --- registro universal de 4 bits: SISO izquierda, SISO derecha, carga paralelo, retención. Controlado por dos entradas de modo (S0, S1).
    \item \textbf{74LS194} --- versión TTL del mismo.
\end{itemize}
\end{notabox}

\subsection{Contadores basados en registros de desplazamiento}

Conectando la salida del último flip-flop a la entrada del primero (registro en anillo), se obtienen secuencias especiales:

\textbf{Contador en anillo:} desplaza un único bit activo de forma circular. Un registro de N etapas genera una secuencia de N estados únicos y requiere N flip-flops para contar N estados (sin decodificador adicional).

\textbf{Contador Johnson (o de anillo torcido):} la salida complementada del último flip-flop ($\overline{Q}$) se realimenta a la entrada del primero. Un registro de N etapas genera una secuencia de 2N estados.

\section{Contadores}

Un contador es un circuito secuencial formado por flip-flops que recorre una secuencia determinada de estados binarios, uno por cada pulso de reloj. El \textbf{módulo} de un contador es la cantidad total de estados distintos que recorre antes de repetirse.

$$\text{Módulo} = 2^N \quad \text{(para un contador binario de N flip-flops)}$$

\subsection{Contadores asíncronos (con propagación)}

El reloj se aplica solo al primer flip-flop. Cada flip-flop siguiente se dispara con la salida del anterior. Los cambios de estado se propagan de etapa en etapa, por eso también se los llama \textbf{contadores con propagación} (\textit{ripple counters}).

\tikzCircuito{cap1/contador-asincrono-4bit}{Contador asíncrono (ripple) de 4 bits con flip-flops JK en modo toggle. La salida Q de cada FF actúa como reloj del siguiente (figura 8.1 Floyd).}

\textbf{Diagrama de tiempos de un contador asíncrono de 3 bits:}

\bannerRevision
\begin{figure}[H]
\centering
% ============================================================
%  Contador asíncrono de 3 bits: diagrama de tiempos
%  - Los FF basculan en el flanco DESCENDENTE de su reloj:
%      Q0 cambia en cada flanco ↓ de CLK   (t = 2, 4, 6, ... 16)
%      Q1 cambia en cada flanco ↓ de Q0    (t = 4, 8, 12, 16)
%      Q2 cambia en cada flanco ↓ de Q1    (t = 8, 16)
%  - Se dibuja ideal (sin retardo de propagación).
%  Coordenadas: ver tikz/estilos.tex (1 carácter = 1 unidad en x).
% ============================================================
\begin{tikztimingtable}[timing/aed, timing/xunit=5mm]
  CLK   & 8{LH}L \\
  $Q_0$ & 2L 3{2H2L} 2H L \\
  $Q_1$ & 4L 4H 4L 4H L \\
  $Q_2$ & 8L 8H L \\
  \extracode
    \begin{pgfonlayer}{background}
      % flancos descendentes del CLK
      \vertlines{2,4,...,16}
    \end{pgfonlayer}
    % conteo en decimal (Q2 Q1 Q0) debajo del diagrama
    \foreach \x/\n in {1/0,3/1,5/2,7/3,9/4,11/5,13/6,15/7,16.5/0}
      \node[font=\sffamily\small] at (\x,{-3*\rowdist-1.0}) {\n};
    \node[font=\sffamily\small, anchor=east] at (-0.8,{-3*\rowdist-1.0}) {Cuenta};
\end{tikztimingtable}

\caption{Diagrama de tiempos del contador asíncrono de 3 bits (módulo 8).
  Q0 es el bit de menor peso (LSB) y Q2 el de mayor peso (MSB).
  Cada salida divide por 2 la frecuencia de la señal anterior.}
\end{figure}

La secuencia de salidas Q2 Q1 Q0 va de 000 a 111 (0 a 7), y luego vuelve a 000. Esto es un contador módulo-8.

\textbf{Desventaja:} el retardo de propagación se acumula etapa a etapa. En contadores muy rápidos o con muchas etapas, los estados intermedios falsos (\textit{glitches}) pueden causar problemas en la lógica de decodificación.

\textbf{Contador de décadas asíncrono:} un contador que llega solo hasta 9 (BCD) y luego se reinicia. Se obtiene forzando el reset del contador cuando se alcanza la combinación binaria 10 ($1010_2$).

\subsection{Contadores síncronos}

El reloj se conecta \textbf{simultáneamente} a todos los flip-flops. Los cambios de estado ocurren al mismo tiempo en todas las etapas, eliminando el problema de acumulación de retardos.

\tikzCircuito{cap1/contador-sincrono-4bit}{Contador síncrono de 4 bits con flip-flops JK (figura 8.16 Floyd)}

La lógica de habilitación de cada etapa se diseña para que el flip-flop cambie de estado solo cuando todos los bits de menor peso estén en 1 (condición de acarreo).

\textbf{Ventajas sobre los asíncronos:}
\begin{itemize}
    \item Todos los bits cambian simultáneamente $\rightarrow$ no hay glitches por propagación.
    \item Pueden operar a mayor frecuencia.
    \item Son preferidos en diseños digitales modernos.
\end{itemize}

\subsection{Contadores ascendentes/descendentes}

Un contador puede recorrer su secuencia en sentido creciente (ascendente, \textit{up}) o decreciente (descendente, \textit{down}). Con una entrada de control de dirección se puede elegir el sentido en tiempo real.

\subsection{Módulo de un contador y truncamiento de secuencia}

Un contador de N flip-flops tiene módulo $2^N$ por defecto. Para obtener un módulo M menor, se detecta el estado M (cuando el contador llega a M) y se fuerza un reset síncrono o asíncrono, haciendo que el contador salte desde M-1 de vuelta a 0.

\textbf{Ejemplo:} contador módulo-5 con 3 flip-flops (que tendrían módulo 8 por defecto). Se detecta el estado 5 ($101_2$) y se usa para resetear el contador, obteniendo la secuencia 0$\rightarrow$1$\rightarrow$2$\rightarrow$3$\rightarrow$4$\rightarrow$0$\rightarrow$\ldots

\subsection{Contadores en cascada}

Se pueden conectar varios contadores en serie para obtener módulos más grandes. La señal de fin de cuenta (\textit{terminal count}, TC) del primer contador habilita el reloj del segundo.

\tikzCircuito{cap1/contadores-cascada}{Dos contadores de 4 bits en cascada formando un contador de 8 bits (figura 8.38 Floyd)}

\subsection{Decodificación de estados}

Para detectar un estado específico del contador se puede usar una puerta AND que active su salida solo cuando las salidas del contador forman la combinación deseada. Se debe tener cuidado con los \textit{glitches} en contadores asíncronos.

\section{Máquinas de estado}

Una máquina de estado finito (FSM, \textit{Finite State Machine}) es un modelo para describir sistemas secuenciales de forma sistemática. Se define por:

\begin{itemize}
    \item Un conjunto de \textbf{estados} posibles.
    \item Un \textbf{estado inicial}.
    \item Las \textbf{condiciones de transición} entre estados (dependientes de las entradas).
    \item Las \textbf{salidas} asociadas a cada estado o a cada transición.
\end{itemize}

\subsection{Diagrama de estados}

El comportamiento de una máquina de estados se representa gráficamente con un \textbf{diagrama de estados}:

\begin{itemize}
    \item Cada estado se representa con un círculo o nodo.
    \item Las transiciones entre estados se representan con flechas etiquetadas con la condición de entrada que las provoca.
    \item Las salidas se indican dentro del nodo (Moore) o sobre la flecha (Mealy).
\end{itemize}

\bannerRevision
\begin{figure}[H]
\centering
\begin{tikzpicture}[
  ->,>=Stealth,
  auto,
  node distance=3.2cm,
  every state/.style={circle, draw=teal!70!black, fill=teal!8,
                      minimum size=1.3cm, font=\sffamily\bfseries},
  initial text={\sffamily\footnotesize Inicio},
  every edge/.style={draw, font=\footnotesize\itshape},
  loop/.style={looseness=6},
]
  % Estados (grilla 2×2)
  \node[state, initial above] (s0)              {$S_0$};
  \node[state]                (s1) [right=of s0]{$S_1$};
  \node[state]                (s2) [below=of s1]{$S_2$};
  \node[state, accepting]     (s3) [below=of s0]{$S_3$};

  % Transiciones
  \path
    (s0) edge[loop left]   node{$x{=}0$}  (s0)
    (s0) edge              node{$x{=}1$}  (s1)
    (s1) edge              node{$x{=}1$}  (s2)
    (s1) edge[bend left]   node{$x{=}0$}  (s0)
    (s2) edge              node{$x{=}0$}  (s3)
    (s2) edge[loop right]  node{$x{=}1$}  (s2)
    (s3) edge[bend left=20] node[left]{$x{=}0/1$} (s0);
\end{tikzpicture}

\caption{Ejemplo de diagrama de estados (FSM de Moore con 4 estados).
  El estado inicial $S_0$ tiene doble flecha de entrada; $S_3$ es el estado
  de aceptación (doble círculo). Las etiquetas sobre las flechas indican la
  condición de entrada que provoca la transición.}
\end{figure}

\subsection{Máquina de Moore}

Las salidas dependen \textbf{únicamente del estado actual}. Son más fáciles de diseñar y menos propensas a glitches en las salidas.

\subsection{Máquina de Mealy}

Las salidas dependen del \textbf{estado actual y de las entradas}. Pueden reaccionar más rápido a los cambios de entrada, pero las salidas pueden tener transitorios.

\subsection{Tabla de estados}

Una forma alternativa y más formal de describir una máquina de estados es mediante una tabla:

\begin{table}[H]
\centering
\begin{tabular}{cccc}
\toprule
Estado actual & Entrada & Estado siguiente & Salida \\
\midrule
S0 & 0 & S0 & 0 \\
S0 & 1 & S1 & 0 \\
S1 & 0 & S0 & 1 \\
S1 & 1 & S2 & 0 \\
\ldots & \ldots & \ldots & \ldots \\
\bottomrule
\end{tabular}
\end{table}

\subsection{Proceso de diseño de una máquina de estados}

\begin{enumerate}
    \item Definir el comportamiento deseado y los estados necesarios.
    \item Elaborar el diagrama de estados.
    \item Asignar una codificación binaria a cada estado.
    \item Derivar las ecuaciones de estado siguiente y de salida (usando mapas de Karnaugh si es necesario).
    \item Implementar el circuito con flip-flops y lógica combinacional.
\end{enumerate}
